\documentclass[11pt]{amsart}

\usepackage[T1]{fontenc}
\usepackage{lmodern}
\usepackage{microtype}
\usepackage{amsmath,amssymb}
\usepackage[hidelinks]{hyperref}

\hypersetup{
  pdftitle={A stabilization counterexample to Conjecture 4.2 of Boix and Lima--Pereira}
}

\newtheorem{theorem}{Theorem}
\newtheorem{proposition}[theorem]{Proposition}
\theoremstyle{remark}
\newtheorem{remark}[theorem]{Remark}

\newcommand{\Tor}{\operatorname{Tor}}
\newcommand{\m}{\mathfrak m}

\title[A stabilization counterexample]
{A Stabilization Counterexample to Conjecture 4.2 of
Boix and Lima--Pereira}

\date{}
\subjclass[2020]{13D02}
\keywords{Tor modules, complete intersections, Koszul homology}

\begin{document}

\begin{abstract}
Boix and Lima--Pereira conjectured a Tor-length identity for complete-intersection
ideals under the restriction $k\leq n-1$. We prove a stabilization result that
preserves the relevant Tor modules and colength, so imposing any fixed gap
$n-k\geq c$ is equivalent to the unrestricted assertion. Stabilizing their
Example~3.1 gives a minimal counterexample with $(n,k)=(4,3)$: the Tor-length
vector is $(4,14,14,4)$, whereas the conjectured vector is $(4,12,12,4)$. The
example is nonreduced and does not address the motivating ICIS conjecture.
\end{abstract}

\maketitle

Let $(R,\m)$ be a regular local ring of dimension $n$, let $I$ be generated by
an $R$-regular sequence of length $k$, and let $J$ be generated by an
$R$-regular sequence of length $n$. Boix and Lima--Pereira conjectured that,
when $k\leq n-1$,
\begin{equation}\label{eq:identity}
 \ell_R\!\left(\Tor_i^R(R/I,R/J)\right)
 =\binom{k}{i}\ell_R\!\left(R/(I+J)\right),
 \qquad 0\leq i\leq k.
\end{equation}
This is \cite[Conjecture~4.2]{BLP}; the unrestricted form $k\leq n$ was
proposed in \cite[Conjecture~4.9]{LPNOT} and disproved in
\cite[Section~3]{BLP} by examples with $k=n$.

For $c\geq0$, let $\mathsf T_c$ denote the universal validity of
\eqref{eq:identity} under the restriction $k\leq n-c$.

\begin{theorem}\label{thm:main}
For every integer $c\geq0$, the assertions $\mathsf T_c$ and $\mathsf T_0$ are
equivalent. Consequently, Conjecture~4.2 of \cite{BLP} is false. Explicitly,
let
\[
 R=\mathbb Q[x,y,z,t]_{(x,y,z,t)},
\]
\[
 I=(x^2,y^2,z^2),\qquad
 J=(x^2+xy,\,xz+z^2,\,y^2+yz,\,t).
\]
Then $I$ and $J$ are generated by regular sequences of lengths $3$ and $4$,
and
\[
 \left(\ell_R\!\left(\Tor_i^R(R/I,R/J)\right)\right)_{i=0}^3
 =(4,14,14,4),
 \qquad
 \ell_R(R/(I+J))=4.
\]
Thus \eqref{eq:identity} predicts $(4,12,12,4)$ and fails for $i=1,2$.
Moreover, $(n,k)=(4,3)$ is minimal in both coordinates.
\end{theorem}

\newpage
\begin{proposition}[Stabilization]\label{prop:stabilization}
For $r\geq1$, set
\[
 R_r=R[t_1,\ldots,t_r]_{\m R[t_1,\ldots,t_r]+(t_1,\ldots,t_r)},
 \qquad I_r=IR_r,
\]
\[
 J_r=JR_r+(t_1,\ldots,t_r).
\]
Then $R_r$ is regular local of dimension $n+r$, and $I_r$ and $J_r$ are
generated by regular sequences of lengths $k$ and $n+r$. For every $i$,
\[
 \Tor_i^{R_r}(R_r/I_r,R_r/J_r)
 \cong \Tor_i^R(R/I,R/J),
 \qquad
 R_r/(I_r+J_r)\cong R/(I+J),
\]
and the corresponding lengths are equal.
\end{proposition}

\begin{proof}
The local map $R\to R_r$ is faithfully flat, so the chosen generators of $I$
and $J$ remain regular after extension. Moreover,
\[
 R_r/JR_r\cong
 (R/J)[t_1,\ldots,t_r]_{(\m/J,t_1,\ldots,t_r)},
\]
and $t_1,\ldots,t_r$ is regular on this ring. This proves the assertions about
regular sequences.

If $f_1,\ldots,f_k$ generate $I$, their Koszul complex is a free resolution of
$R/I$. Since $R_r/J_r\cong R/J$, there is an isomorphism of complexes
\[
 K_{R_r}(f_1,\ldots,f_k)\otimes_{R_r}R_r/J_r
 \cong K_R(f_1,\ldots,f_k)\otimes_R R/J.
\]
Taking homology gives the Tor isomorphisms, and the quotient isomorphism is
immediate. The modules on the left are annihilated by $(t_1,\ldots,t_r)$, so
length over $R_r$ equals length over
$R_r/(t_1,\ldots,t_r)\cong R$.
\end{proof}

\begin{proof}[Proof of Theorem~\ref{thm:main}]
The case $c=0$ is tautological. Let $c\geq1$. The implication
$\mathsf T_0\Rightarrow\mathsf T_c$ is immediate. Conversely, assume
$\mathsf T_c$ and take an arbitrary instance of $\mathsf T_0$ in dimension
$n$. Stabilizing with $r=c$ gives dimension $n+c$ and
\[
 k\leq n=(n+c)-c.
\]
Thus $\mathsf T_c$ applies, and Proposition~\ref{prop:stabilization} transports
\eqref{eq:identity} back to the original instance.

We now verify the displayed counterexample exactly. It is the one-variable
stabilization of \cite[Example~3.1]{BLP}. Set
\[
 S=\mathbb Q[x,y,z],\qquad
 J_0=(x^2+xy,\,xz+z^2,\,y^2+yz),\qquad A=S/J_0.
\]
For graded reverse lexicographic order with $x>y>z$, a Gr\"obner basis of
$J_0$ is
\[
 x^2+xy,\quad xz+z^2,\quad y^2+yz,\quad yz^2-z^3,\quad z^4.
\]
Hence $A$ has homogeneous basis
\[
 1;\qquad x,y,z;\qquad xy,yz,z^2;\qquad z^3.
\]
In particular, $(x,y,z)^4\subseteq J_0$, so $J_0$ is $(x,y,z)$-primary and its
three generators form a regular sequence in $S_{(x,y,z)}$. The sequence
$x^2,y^2,z^2$ is regular there as well.

In $A$, put $q_1=x^2$, $q_2=y^2$, and $q_3=z^2$. The defining relations give
\[
 q_1=-xy,\qquad q_2=-yz,\qquad q_3=z^2,
\]
so $q_1,q_2,q_3$ form a basis of $A_2$. Relative to this basis, the basis
$x,y,z$ of $A_1$, and the basis $z^3$ of $A_3$, multiplication
$A_2\otimes A_1\to A_3$ is represented by
\[
 T=\begin{pmatrix}
 1&-1& 1\\
 1& 1&-1\\
-1& 1& 1
 \end{pmatrix},
 \qquad \det T=4.
\]
Thus this multiplication pairing is perfect.

Consider the Koszul complex $K(q_1,q_2,q_3;A)$. Its chain dimensions are
$8,24,24,8$. Since the $q_i$ have degree $2$ and $A_d=0$ for $d\geq4$, only
coefficient degrees $0$ and $1$ contribute to its differentials. Let
$V=\mathbb Q^3$ with basis corresponding to the $q_i$. The two nonzero pieces
of $d_2$ identify with
\[
 \bigwedge^2V\longrightarrow V\otimes V,
 \qquad u\wedge v\longmapsto u\otimes v-v\otimes u,
\]
and, using the perfect pairing above, with
\[
 V^*\otimes\bigwedge^2V\longrightarrow V,
 \qquad
 \varphi\otimes(u\wedge v)\longmapsto
 \varphi(u)v-\varphi(v)u.
\]
The first map is injective and the second is surjective, so both have rank $3$.
The degree-zero and degree-one parts of $d_1$ have ranks $3$ and $1$.
Those of $d_3$ have ranks $1$ and $3$, the latter map being an isomorphism
under the same perfect pairing. Hence
\[
 \operatorname{rank}d_1=4,\qquad
 \operatorname{rank}d_2=6,\qquad
 \operatorname{rank}d_3=4.
\]
Therefore the homology dimensions are
\[
 (8-4,\ 24-4-6,\ 24-6-4,\ 8-4)=(4,14,14,4).
\]
Since $A$ is an Artinian local $\mathbb Q$-algebra with residue field
$\mathbb Q$, these dimensions are the corresponding lengths. Also,
$A/(q_1,q_2,q_3)$ has basis $1,x,y,z$ (the $q_i$ span $A_2$ and
$z^3=yq_3$), and hence length $4$.
Proposition~\ref{prop:stabilization} with $r=1$ gives the stated
four-dimensional example.

Finally, the identity holds for $k=2$ by \cite[Theorem~0.2(ii)]{BLP} and for
$k=1$ by \cite[Remark~4.3(ii)]{BLP}. Every counterexample to Conjecture~4.2
therefore has $k\geq3$ and $n\geq k+1\geq4$, proving minimality.
\end{proof}

\begin{remark}[Scope]
The complete intersection $R/I$ is nonreduced, so the example lies outside the
reduced ICIS setting of \cite[Conjecture~4.1]{BLP}. It also satisfies the
inequalities proposed in \cite[Question~5.1]{BLP}, since
$(4,14,14,4)\geq(4,12,12,4)$ componentwise.
\end{remark}

\enlargethispage{2\baselineskip}

\begin{thebibliography}{9}

\bibitem{BLP}
A.~F. Boix and B.~K. Lima--Pereira,
\emph{Exploring a local variant of the Buchsbaum--Eisenbud--Horrocks
conjecture},
Proc. Roy. Soc. Edinburgh Sect. A, First View (2026), 1--20,
\href{https://doi.org/10.1017/prm.2026.10176}
{doi:10.1017/prm.2026.10176}.

\bibitem{LPNOT}
B.~K. Lima--Pereira, J.~J. Nu\~no--Ballesteros,
B.~Or\'efice--Okamoto, and J.~N. Tomazella,
\emph{The Bruce--Roberts numbers of a function on an ICIS},
Q. J. Math. \textbf{75} (2024), no.~1, 31--50,
\href{https://doi.org/10.1093/qmath/haad038}
{doi:10.1093/qmath/haad038}.

\end{thebibliography}

\end{document}